\chapter{Implementation}\label{ch:impl}

This chapter describes the SystemVerilog implementation in \texttt{rtl/} (top level
\texttt{strata\_top}, core \texttt{strata\_core}, Transfer Unit \texttt{strata\_tu}, far bus
\texttt{strata\_bus}, memory targets \texttt{strata\_mem}, devices \texttt{strata\_dev}). It is not
normative. Figures marked \estimate are projections.

\section{Organisation}

\begin{figure}[htbp]
\centering
\begin{tikzpicture}[b/.style={draw=accent,fill=accentlight,rounded corners=2pt,minimum height=8mm,
    align=center,font=\small},t/.style={draw=black!60,fill=black!5,minimum height=8mm,align=center,font=\small},
    >={Stealth[length=2mm]}]
  \foreach \i/\x in {0/0,1/3.4,2/6.8} {
    \node[b,minimum width=28mm] (core\i) at (\x,0) {core \ifnum\i=2 $N\!-\!1$\else\i\fi\\\scriptsize I-LS, D-LS (dual-port RAM)};
    \node[b,minimum width=13mm,font=\scriptsize] (tu\i) at (\x-0.75,-1.3) {TU};
    \node[b,minimum width=13mm,font=\scriptsize] (am\i) at (\x+0.75,-1.3) {atomics};
    \draw[<->] (core\i.south -| tu\i) -- (tu\i); \draw[<->] (core\i.south -| am\i) -- (am\i);
  }
  \node[b,minimum width=22mm] (dma) at (10.2,-1.3) {block device\\\scriptsize DMA engine};
  \node[draw=black!60,minimum width=118mm,minimum height=7mm,font=\small] (bus) at (5.0,-2.8)
    {far bus: one burst transaction at a time, system wide};
  \foreach \n in {tu0,am0,tu1,am1,tu2,am2,dma} \draw[<->] (\n.south) -- (\n.south |- bus.north);
  \node[t,minimum width=30mm] (t3) at (0.5,-4.3) {T3 system memory};
  \node[t,minimum width=30mm] (t2) at (4.6,-4.3) {T2 cluster store};
  \node[t,minimum width=36mm] (dev) at (9.0,-4.3) {device space\\\scriptsize UART, platform ROM, block device};
  \foreach \n in {t3,t2,dev} \draw[<->] (\n.north) -- (\n.north |- bus.south);
  \node[draw=accent,dashed,minimum width=118mm,minimum height=6mm,font=\small] (ib) at (5.0,1.3) {capability inboxes: per-sender queues, 2 deep, one RAM per core; nodes $0 \ldots N$};
  \foreach \i in {0,1,2} \draw[<->,dashed,accent] (core\i.north) -- (core\i.north |- ib.south);
  \draw[<->,dashed,accent] (dma.north) -- (dma.north |- ib.south);
\end{tikzpicture}
\caption{Structure of the implementation}\label{fig:impl}
\end{figure}

\subsection{Core}
The core is multi-cycle over block-RAM Local Stores. Its states are: RST, after reset, clearing the
integer registers and the capability save table (one register and one slot per cycle, 256 cycles);
F0, present the fetch address to the I-LS (stalling while an in-flight \insn{tgetx} covers it); F1,
latch the instruction; EX, execute, remaining there while stalled; CT, the second cycle of the
save-table instructions; LD, receive load data; LG, LT and LW, the set-entry computation and the
tag-table walk of \insn{ltag}/\insn{linst}; DV, multiplication (four partial products of a
$32\times32$ multiplier, one per cycle) and iterative division (one quotient bit per cycle); AM, a far
atomic in flight. Each Local Store is a true dual-port RAM: port A belongs to the core and port B to
its TU. The integer registers are LUT-RAM with one write port and three asynchronous read ports (every
result is written through the one port at the end of its cycle); the capability registers, the
sixteen-entry overlay map, the mode bits and the banked trap CSRs are flip-flops. A pending stop
request is checked in F0 and in EX, so a user-mode core takes it at the next instruction boundary
even while an instruction is stalled.

Every capability satisfies $\mathit{base} + \mathit{length} \le$ the size of its space, however it was
made, so the implementation holds a capability's base and length in $\mathit{FW} = 27$ bits
($\log_2$ of the largest space, system memory, plus one); the architectural values are these,
zero-extended. A 64-bit register compared with a length is larger when any bit above the narrow part
is set, and otherwise compared in 27 bits; trap values are always the full registers. Two read ports
serve the capability registers: A, the operand (\reg{cs} of a capability instruction, the transfer's,
atomic's or overlay's capability otherwise), and D, the destination (and \insn{csend}'s endpoint). Each
cycle makes at most one capability write to \reg{cd} and one to \reg{cs}. Likewise an overlay map
entry holds its far address in 27 bits and its I-LS address in 16, and its range match is one 27-bit
subtraction; the target address is computed once, for the matching entry, in 64 bits. Local Store
addresses and lengths (in the TU, the hazard checks and the range checks) are 17 bits. pc$+4$ and
branch and \insn{jal} targets are computed on 17 and 24 bits and sign-extended, which is exact (the pc
is inside the I-LS and offsets are below $2^{23}$), so a target outside the I-LS faults with its full
address. Loads, stores and \insn{jalr} take their address from the ALU's adder, which also gives the
branch conditions.

The capability save table ($\mathit{CTAB}$ = 256 entries, a parameter of \texttt{strata\_pkg}) is a
block RAM with a registered read, one capability per word (an empty slot holds the null
capability). EX reads the slot of \insn{csave}, \insn{crestore}, \insn{csavex} or \insn{crestorex}
(the instruction's immediate or, for the register forms, the low eight bits of \reg{rx}, a larger
value trapping first), and the instruction finishes in CT with the entry, one cycle later.
\insn{ltag} and \insn{linst} compute their set's entry address
$\mathit{ttag} + \mathit{set}\times(\mathit{ways}+1)\times 8$ in 64 bits (it is the \cause{lsfault}
value) by shift-and-add in the first cycle of LG and check it in the second. The decoder is purely
combinational and flags undefined encodings.

Table~\ref{tab:cycles} gives the cycles an instruction takes when nothing stalls it.

\begin{table}[htbp]
\centering\small
\caption{Cycles per instruction, without stalls}\label{tab:cycles}
\begin{tabular}{@{}l r@{}}
\toprule
\textbf{Instructions} & \textbf{Cycles}\\\midrule
most (ALU, branches, stores, capability, CSR, transfer issue) & 3\\
loads & 4\\
\insn{csave}, \insn{crestore}, \insn{csavex}, \insn{crestorex} & 4\\
\insn{mul} & 6\\
\insn{mulhu} & 7\\
division and remainder & about 68\\
\insn{ltag}, \insn{linst} & $5 + 2k$ (\insn{ltag}: $k$ tag words read) or 9 (\insn{linst})\\
\bottomrule
\end{tabular}
\end{table}

\subsection{Transfer Unit}
The TU accepts transfer commands from EX in program order into a four-entry in-order queue and executes
them one at a time as far-bus bursts, so one core's transfers take effect at the far side in issue
order (invariant I2). Issue waits while the queue is full, while the transfer's tag is busy, or while
the D-LS range of a \insn{tget} or \insn{tput} overlaps a queued device read (with CCR.PTRAP the core
traps with \cause{pend} instead). For every
queued entry the TU compares its Local Store range with the range of the current load, store or
tag-table access and with the fetch address, and raises the stalls that implement invariant I1 (an
entry keeps its range's end, so these are comparisons, not additions).
The TU reports two status signals to the core: \emph{idle} (its queue is empty) and \emph{memory
idle} (no queued or executing transfer targets a space other than device space). \insn{cgive},
\insn{csend}, \insn{cshare}, \insn{catomic} and \insn{tamo} wait for memory idle, the publication rule
of invariant I4; device transfers, such as a pending console break read, do not hold them up.
\insn{halt}, in either mode, waits for idle, device reads included. Memory
transfers move 8-byte beats; device transfers move single bytes to consecutive device offsets.

\subsection{Far bus and targets}
The far bus connects the masters (every core's TU and atomic unit, and the block device's DMA engine) to
the targets (T3, T2, device space). Exactly one burst transaction is in flight system-wide, which keeps
each core's transfers in order at the target and makes atomics trivially linearizable. A memory target
has a fixed start-up latency per burst (by default 20 cycles for T3 and 4 for T2) and then delivers one
beat per cycle; an atomic is a read-modify-write of one word. A device with nothing to deliver yet
(console input with no byte pending, the break port, WAIT before the deadline, VSYNC before the next
refresh) answers \emph{retry} before any data, so a pending device read does not hold the bus: the TU
tries it again later, and the rest of the machine runs meanwhile.

\subsection{Revocation granules and core control}
The granule table (a 16-bit generation and a head bit for each of the 256 granules of T3 and the one of
T2, Section~\ref{sec:domains}) exists once in the top level and once in every core. Each core's copy is
LUT-RAM with two asynchronous read ports, so its \texttt{Live} checks stay combinational: port A
serves the operand capability (\reg{cs} of a capability instruction, the transfer's, atomic's or
overlay's capability otherwise; never both in one instruction) and port D the destination register.
In CT the save-table slot's entry takes the port its instruction does not need (\insn{csave}: D,
\insn{crestore}: A), so all three checks are still made in the same cycle. Only the top level writes the table: a core executing \insn{cdom},
\insn{crevoke} or \insn{cundom} raises a request and stalls in EX; the top level serves the
lowest-numbered requester with a walker that reads its master copy and writes one granule per cycle,
broadcasting every write to all the copies (for \insn{crevoke}: back to the head of the first touched
unit, a read pass to the end of the last, a write pass), and then signals the core, which writes its
capability registers and retires. A walk costs about two cycles per granule of the touched units; other
cores keep running meanwhile. The alternative, one shared table behind a pipeline stage, would add an
arbitrated port to every capability check of every core; three small read ports per core (256 entries
of 17 bits) are cheaper. Revocation by address cost about 1,000 LUT sites fewer than the
fifteen-domain table it replaced, which held a full region register per domain. Core-control space is
a far-bus target in the top level: a \insn{tput} through a core-control capability raises the stop request of the core whose doubleword it
writes with bit~0 set.

\subsection{Devices}
Device space contains the UART console, the platform table ROM, block device 0, the timer and the
framebuffer controller. A read of the console break port answers \emph{retry} until a Ctrl-C is counted
or the release port is written while it waits, so a pending break read occupies its tag but never the
bus. The device keeps one waiting bit and one released bit per core, so that one write to the release
port completes the break reads of all cores. The timer is a 64-bit counter of clock cycles divided by
the parameter \texttt{CLK\_MHZ}; in simulation the harness can supply host time instead. NOW and
DEADLINE are latched at byte 0 of a read and DEADLINE is committed at byte 7 of a write. The
controller counts refreshes of 16667 ticks for VSYNC; PRESENT calls the simulation harness, which
shows the framebuffer, and has no effect in synthesis. The framebuffer itself is T3 memory: the
harness reads back the range written since its last repaint. On FPGA builds a received byte 3 counts as
Ctrl-C and is not delivered to the input port; the simulation harness reports Ctrl-C from the host
terminal or from \texttt{-b}. The block device
receives its buffer capability through its own inbox (node $N$), moves blocks between its storage and the
memory the capability names using its own bus port, posts a completion and gives the capability back.
Holding the capability is the only way it can reach that memory.

\section{Toward FPGA}

The RTL is synthesizable. Its FPGA-specific elements are:
\begin{itemize}
\item Local Stores in inferred dual-port block RAM, initialised with \texttt{\$readmemh} for the boot ROM;
\item an iterative divider in place of a combinational one;
\item an 8N1 UART console (115200 baud at 100~MHz) with a 16-byte receive FIFO;
\item an external system-memory port: with the top-level parameter \texttt{T3\_EXT} set, T3 is not built
  in block RAM but exported as a target port. \texttt{fpga/rtl/strata\_t3\_litedram.sv} adapts that port
  to a LiteDRAM native user port; each 8-byte beat becomes one DRAM access, and an atomic is a read
  followed by a write while the far bus is held, so atomics stay linearizable.
\end{itemize}

\texttt{fpga/arty/} contains a top level for the Digilent Arty A7-100T: two cores, a 64~KiB cluster
store and a 128~KiB disk in block RAM, the console on the USB UART, and T3 in the board's DDR3 through a
LiteDRAM controller. It is synthesizable but has not yet run on hardware, and as the next section shows
two cores are more than that board holds. \texttt{fpga/synth.sh}
synthesizes the system with Yosys (through \texttt{sv2v}) for Xilinx 7-series or Lattice ECP5.

\subsection{Synthesis results}

Table~\ref{tab:synth} gives the Yosys results (\texttt{synth\_xilinx} and \texttt{synth\_ecp5}, Yosys
0.70) for systems of one, two and six cores with a 64~KiB cluster store, a 128~KiB block-RAM disk and a
32~KiB block-RAM stand-in for system memory.

\begin{table}[htbp]
\centering\small
\caption{Synthesis results (Yosys, before placement)}\label{tab:synth}
\begin{tabular}{@{}l r r r r@{}}
\toprule
\textbf{System} & \textbf{Logic} & \textbf{Flip-flops} & \textbf{Block RAM} & \textbf{Multipliers}\\\midrule
7-series, 1 core & 30,530 LUTs & 6,871 & 88 RAMB36 + 4 RAMB18 & 4 DSP48E1\\
7-series, 2 cores & 58,066 LUTs & 11,713 & 120 RAMB36 + 7 RAMB18 & 8 DSP48E1\\
7-series, 6 cores & 172,984 LUTs & 31,104 & 248 RAMB36 + 19 RAMB18 & 24 DSP48E1\\
ECP5, 1 core & 42,608 LUT4 sites & 6,831 & 180 DP16KD & 4 MULT18X18D\\
\bottomrule
\end{tabular}
\end{table}

These are pre-placement figures: nothing has been placed and routed, and no clock frequency is known.
The ECP5 column counts LUT4 sites, that is LUT4 cells plus two for each carry cell. One more core costs
about 27,500 LUTs, 4,800 flip-flops, 32 RAMB36 (its two 64~KiB Local Stores), 3 RAMB18 (its save table)
and 4 DSP slices on 7-series. No latches are inferred, and every large memory is block RAM on both
families. The Local Stores are written through both ports, and the ECP5's block RAM does not define what
one port reads from a byte that the other is writing; the hazard stalls of invariant I1 guarantee that
nobody reads such a byte, and the RAM is marked accordingly so that it maps to block RAM on both
families.

An earlier version of the core took about 47,000 LUTs (50,606 for the one-core system, 80,999 LUT4
sites on the ECP5, 290,431 LUTs for six cores). Most of the difference is the narrow fields of the
previous sections (the overlay map alone was about 12,000 LUTs, its entries holding 64-bit addresses),
fewer register-file ports, and RAM for the save table, the integer registers and the inboxes;
\texttt{fpga/README.md} has the steps and their effect.

\begin{table}[htbp]
\centering\small
\caption{FPGA boards and what fits them, from the synthesis results}\label{tab:fpga}
\begin{tabularx}{\linewidth}{@{}l X X X@{}}
\toprule
 & \textbf{Digilent Arty A7-100T} & \textbf{Radiona ULX3S (ECP5-85F)} & \textbf{Digilent Nexys Video}\\\midrule
Device & XC7A100T & LFE5U-85F & XC7A200T\\
Logic & 63,400 LUTs & 83,640 LUT4s & 134,600 LUTs\\
On-chip RAM & 4.86 Mb BRAM (135 $\times$ 36 Kb) & $\approx$3.7 Mb EBR (208 $\times$ 18 Kb) & 13.1 Mb BRAM (365 $\times$ 36 Kb)\\
External memory & 256 MB DDR3 & 32 MB SDRAM & 512 MB DDR3\\
Toolchain & Vivado or yosys & open: yosys + nextpnr & Vivado or yosys\\
Cores & 1 (48\% of the LUTs) & 1 (51\% of the LUT4s; not placed and routed) & 2 (43\% of the LUTs)\\
T2 & 64 KiB & 64 KiB & 64 KiB\\
T3 & DDR3 via LiteDRAM (in RTL) & SDRAM (no adapter yet) & DDR3 via LiteDRAM\\
\bottomrule
\end{tabularx}
\end{table}

Each core also needs 1~Mb of block RAM for its Local Stores (64 DP16KD), which on the ECP5-85F, with
the cluster store and the disk in block RAM too, limits a system to one core whatever the logic. StrataOS needs six cores, and the 6-core system's 172,984 LUTs are more than
any Artix-7 holds but fit a Kintex-7 325T. The systems of one and two cores run bare-metal programs.
The FPGA builds have no video output: the framebuffer window is memory that nothing scans out, and the
console is the UART.

\subsection{Schematics}

Figures~\ref{fig:sch-div}--\ref{fig:sch-sys} are netlist schematics generated from the RTL
(\texttt{doc/figures/}). They are included for orientation; zoom into the PDF to read them.

\begin{figure}[p]
\centering
\includegraphics[width=\linewidth]{figures/schematic-divider.pdf}
\caption{Netlist schematic of the iterative divider (\texttt{strata\_div})}\label{fig:sch-div}
\vspace{2em}
\includegraphics[width=\linewidth]{figures/schematic-uart.pdf}
\caption{Netlist schematic of the UART (\texttt{strata\_uart})}\label{fig:sch-uart}
\end{figure}

\begin{figure}[p]
\centering
\includegraphics[height=0.88\textheight]{figures/schematic-system.pdf}
\caption{Netlist schematic of the system top level (\texttt{strata\_top})}\label{fig:sch-sys}
\end{figure}

\section{Conformance}\label{sec:rtl-deltas}

The implementation conforms when \texttt{make test} passes: every example produces the same output,
halt states and per-core instruction traces as the reference emulator; a complete StrataOS boot and
scripted session produces the same output (Section~\ref{sec:compare}); and every test of the
conformance suite passes with the reference emulator's signature (Section~\ref{sec:conformance}). The
RTL implements every feature of this edition, including supervisor and user modes with banked trap
state, \insn{svc} and the user-mode inbox rule, revocation by address, endpoints, core control,
\insn{tany}, pending device reads and their ordering rules, the 8-byte alignment of capabilities, the
timer and the framebuffer controller, sixteen-entry overlay mapping with \csr{xslot}, far links and
\insn{xret}, per-sender inbox queues, \csr{takefrom}, the no-overwrite rule, and the publication rule.
Any remaining difference from this manual is a defect of the RTL.
